\chapter{Att kontrollera ett påstående}
\label{app:verifiera}

Den här modulen sätter namn på tre principer --- DRY, SRP och KISS --- och
att sätta namn på en princip är att göra två faktapåståenden på en gång.
Det ena är ett \emph{ursprung}: principen härrör från den och den.  Det
andra är en \emph{etablering}: den är en spridd och accepterad praktik.
De två kräver olika slags källor, och de kan mycket väl få olika svar.

Att skilja dem åt är hela poängen med den här bilagan.  Ursprunget kan
avgöras av en enda källa --- den där idén först formulerades --- medan
etableringen kräver att man ser efter hur fältet faktiskt använder
principen, och lika mycket vad fältet har invänt mot den.  En källa som
säger att principen är bra bevisar bara att någon tycker det.

\Cref{tab:principerna} sammanfattar var vart och ett av modulens
påståenden är kontrollerat.  \Cref{app:srp,app:kiss} visar arbetsgången
genomförd på SRP respektive KISS och kan läsas fristående.

\section{Gör påståendet till en fråga}

Ett påstående går att kontrollera först när det är tydligt vad som skulle
kunna visa att det är \emph{fel}.  \enquote{Härrör SRP från Robert C.
Martin?} kan besvaras med ja eller nej.  \enquote{SRP är en bra princip}
kan inte det: ingen tänkbar källa skulle visa att det är fel, eftersom
\enquote{bra} är ett omdöme.  Den som säger så menar oftast något som
\emph{kan} prövas --- \enquote{leder SRP till kod som är lättare att
underhålla?} --- och den frågan kan få svaret nej.

\section{Välj rätt sorts källa}

\begin{description}
  \item[Primärkällan] för var en idé \emph{kommer ifrån}.  Vill man veta
    om SRP härrör från Martin läser man Martins egen text, inte en
    lärobok som återger den.  Finns flera utgåvor väljer man den
    tidigaste där idén faktiskt står --- och det är värt att kontrollera,
    för den allmänt antagna förstakällan behöver inte vara rätt (se
    \cref{app:srp}).
  \item[Enskilda studier] för ett specifikt resultat, med förbehållet att
    en enskild studie visar att något observerats en gång.
  \item[Översikter] (\foreignlanguage{english}{reviews}) för om något är
    \emph{etablerat} eller vad forskningen \emph{sammantaget} säger.  En
    systematisk översikt går igenom hundratals studier efter en angiven
    metod.  Att en sådan översikt \emph{saknas} för ett begrepp är i sig
    en upplysning.
  \item[Kurs- och kursplaner, läroböcker] för om något är etablerat i
    \emph{undervisningen}, vilket är en annan fråga än om det är
    etablerat i forskningen.
\end{description}
Var man hittar dem: universitetsbibliotekets söktjänst, och databaser som
Scopus, Web of Science, IEEE Xplore, ACM Digital Library, DBLP (datalogi)
och OpenAlex (öppen; \enquote{öppen} står här för öppna data, inte för
\foreignlanguage{english}{open access}).  I bilagorna användes
kommandoradsverktyget \texttt{scholar}\footnote{Paketet
  \texttt{scholarcli} på PyPI: \url{https://pypi.org/project/scholarcli/}.},
som ställer samma fråga till flera databaser på en gång och sparar allt
som hände i en \emph{session}.  Det man får ut av bibliotekets söktjänst
är detsamma, bara långsammare.

\section{Sök så att du kan ha fel}

Den som söker på namn hon redan känner till hittar det hon väntar sig.
Sök därför \emph{ämnesdrivet} --- på begrepp, inte författare.  Alla
sökningar i bilagorna är gjorda så; källor som bara går att nå som kända
dokument (böcker, artiklar utan \textsc{doi}, dokument på en organisations
egen webbplats) anges som sådana, med skälet.

Några praktiska regler:
\begin{itemize}
  \item Håll frågorna korta.  Flera korta frågor slår en lång.
  \item Ställ \emph{samma} frågor i alla databaser.  Behöver en databas
    en annan syntax översätter man frågan; man byter inte ord.
  \item Sök begreppets alla namn.  Praktikerns ord är sällan det ord
    forskningen indexerar under: litteraturen om SRP-brott står under
    \enquote{god class} och \enquote{design smell}, inte under
    \enquote{SRP violation}.  Det var just den översättningen som fällde
    avgörandet i \cref{app:srp}.
  \item Sök efter \emph{invändningar} med samma begrepp som stödsökningen
    använde.  Ett påstående som överlevt en motsökning är värt mer än ett
    som aldrig prövats, och en invändning som hittas blir ett
    \emph{förbehåll} i svaret, inte ett nederlag.
  \item Blev en sökning dålig --- för bred, fel ämne, fel ordval --- gör
    om den, och behåll allt den gamla hittade.  En omgjord sökning ska
    inte tappa något.
\end{itemize}

\section{Läs i fulltext, och skriv ner hur}

En träff i en databas är inte ett belägg.  Sammanfattningen
(\foreignlanguage{english}{abstract}) räcker inte heller: den säger vad
författarna vill ha sagt, inte vad de visat.  Läs hela texten och leta
upp \emph{det ställe} som styrker påståendet --- ett ordagrant citat med
sidnummer.  Går texten inte att komma åt säger man det, och nöjer sig
inte med att låtsas ha läst den; i \cref{app:srp} är tre källor lästa
bara i sammanfattning, och det står utskrivet vilka.

Skriv sedan ner hur ni gick till väga, så att någon annan kan följa
spåret.  I det här materialet står den anteckningen som en kommentar
ovanför varje post i litteraturlistans källfil, med fasta fält:
\begin{description}
  \item[\texttt{CLAIM}] vilket påstående källan ska styrka;
  \item[\texttt{FOUND-VIA}] hur den hittades --- sökfråga och databas,
    eller \enquote{känt dokument};
  \item[\texttt{PICKED}] varför just den, framför andra träffar;
  \item[\texttt{QUOTE}] det ordagranna citatet, med sida;
  \item[\texttt{VERIFIED}] att fulltexten lästs, och var;
  \item[\texttt{COUNTER}] vad motsökningen gav;
  \item[\texttt{DATE}] när.
\end{description}

\section{Dra slutsatsen --- med förbehåll}

Svaret skrivs ut tillsammans med det som talar emot.  Båda bilagorna
svarar \enquote{ja, med förbehåll}, men förbehållen är av olika slag.
För SRP är ursprunget entydigt men innebörden omtvistad: vad som räknas
som \emph{ett} ansvar visar sig vara en bedömningsfråga som mätbart
skiljer bedömare åt.  För KISS är det tvärtom: att enkelhet är ett
etablerat designvärde går att belägga, medan den vanliga
upphovsuppgiften inte går att bekräfta i någon förstahandskälla.

Förbehållen är inte svagheter i svaren.  De \emph{är} svaren, och de
avgör hur principerna används i materialet: som frågor att ställa om ett
program, inte som regler som avgör saken.

\begin{table}[htbp]
  \centering
  \caption{Modulens påståenden, och var vart och ett är kontrollerat.}
  \label{tab:principerna}
  \footnotesize
  \begin{tabular}{@{}p{0.42\linewidth}p{0.5\linewidth}@{}}
    \toprule
    Påstående & Kontrollerat \\
    \midrule
    Att inte upprepa sig i ett program är en bra princip, och den
      härrör från Hunt och Thomas (DRY)
      & Redan belagt i \emph{Algoritmiskt tänkande}, bilaga E; posten är
        kopierad därifrån med sin proveniensanteckning \\
    SRP härrör från Robert C. Martin, och är en etablerad praktik
      & \Cref{app:srp}: tvåsidig sökning, sessionen för SRP \\
    KISS tillskrivs Kelly Johnson, och enkelhet är ett etablerat
      designvärde
      & \Cref{app:kiss}: tvåsidig sökning, sessionen för KISS \\
    En funktion som skapas men aldrig anropas gör att arbetet uteblir
      & Återanvänt från forskningsartikeln om missuppfattningar \\
    Nybörjare utelämnar returvärdet, och tror att en utskrift sist i
      kroppen fungerar som en return-sats
      & Återanvänt från forskningsartikeln om missuppfattningar \\
    Nybörjare urskiljer sällan att \mintinline{python}{return}
      \emph{avslutar} funktionen
      & Återanvänt från forskningsartikeln om missuppfattningar \\
    Studenter bygger egna regler för lokala och globala variabler
      & Återanvänt från forskningsartikeln om missuppfattningar \\
    PEP~257 är konventionen för docstrings; PEP~8 är stilguiden
      & Primärkällor, hämtade och lästa \\
    Kontrast måste upplevas samtidigt, och komma före generalisering
      & Primärkälla, Martons \emph{Necessary Conditions of Learning} \\
    \bottomrule
  \end{tabular}
\end{table}

\section{En anmärkning om språkmodeller}

Sökningarna i bilagorna gav över tusen träffar var.  En språkmodell
användes för att \emph{sortera} dem --- hör träffen till ämnet över huvud
taget, och åt vilket håll pekar den? --- så att gallringen kunde granskas.
Lägg märke till arbetsfördelningen: modellen sorterade, men varje källa
som citeras är läst och vald av en människa, och även sorteringen
granskades av en människa där modellen var osäker.  Modellens klassning
står därför med sin konfidens i tabellerna, så att den kan ifrågasättas.
Använd gärna språkmodeller för att hitta och sortera; låt dem aldrig vara
det sista ledet.  De kan göra fel, men det är ni själva som måste stå
till svars för innehållet.

\chapter{Är SRP Robert C. Martins princip, och är den etablerad?}
\label{app:srp}
\chapterprecis{Författaren har ännu inte granskat resultaten i den här
  bilagan i sin helhet.}

I \cref{sec:dela-upp} säger vi att principen om ett enda ansvar --- en
modul ska aldrig ha mer än ett skäl att ändras --- är Robert C. Martins,
och att den är en etablerad praktik.  Det är två faktapåståenden av olika
slag, så frågan är: \emph{härrör SRP från Martin, och i vilken mening är
den etablerad?}  \Cref{sec:srp-metod} beskriver sökningen,
\cref{sec:srp-resultat} vad den fann, \cref{sec:srp-diskussion} vad som
begränsar svaret och \cref{sec:srp-slutsats} ger det.

\section{Metod}
\label{sec:srp-metod}

Frågan söktes ämnesdrivet och tvåsidigt den 3 september 2026 med
verktyget \texttt{scholar} mot de databaser som svarade: OpenAlex, DBLP,
Web of Science, IEEE Xplore och Scopus.  Semantic Scholar avvisade sin
nyckel och saknas därför helt; varje sådan cell är en databas som var
nere, inte ett nollresultat.  Stödfrågorna söker principen under sina
egna namn --- \emph{single responsibility principle}, \emph{SOLID} ---
och under de begrepp den bygger på: kohesion, ansvarsdriven design,
informationsdöljande.  Motfrågorna söker med samma begrepp efter kritik,
och lägger därtill till fältets \emph{indexerade} namn för ett brott mot
principen: \emph{god class}, \emph{large class}, \emph{blob},
\emph{design smell}.  Den översättningen är avgörande.  Praktikerns
uttryck \enquote{SRP violation} förekommer knappt i litteraturen, medan
\enquote{God Class} är det ord den empiriska forskningen faktiskt
använder --- och det var den motfrågan som gav kapitlets starkaste
invändningar.  \Cref{tab:sok-queries-srp} visar frågorna och antalet
träffar per databas.

Två stödfrågor gjordes om och de ursprungliga är strukna.  Båda använde
det oavgränsade uttrycket \enquote{design principles}, som inte är
datalogiskt: träfflistorna fylldes av kemi, materialvetenskap, medicin
och utbildningsforskning utanför datalogin.  De omgjorda frågorna binder
uttrycket till SOLID, SRP eller objektorienterad design, och varje
datalogiskt relevant träff från de gamla frågorna finns kvar i de nya.
En lucka upptäcktes därtill och stängdes genom en bättre begreppsfråga,
inte genom att leta efter den saknade artikeln vid namn: fråga~S8 parar
principens namn med ord för omstrukturering och modularitet, och fångar
därmed det arbete om automatiserad uppdelning av långa metoder som de
tidigare frågorna missade.

DBLP:s nollor kräver en förklaring.  DBLP indexerar bara titlar och
struntar i citattecken och booleska operatorer, så en flerledad
begreppsfråga matchar nästan ingenting.  De frågorna ställdes därför en
gång till i kort, titelformad gestalt (raderna märkta \textsuperscript{b}).
En nolla även där betyder \enquote{ingen sådan titel i DBLP}, inte
\enquote{ingen sådan litteratur}.

Två källor hämtades som kända dokument, av skäl som är databasernas.
Martins bok indexeras inte som artikel, så kapitlet lästes i författarens
egen kapitelutskrift via Internetarkivet och bokens uppgifter hämtades
från förlagets katalogpost.  Martins tidigare artikel från 2000 är grå
litteratur som ingen av databaserna indexerar och hämtades på samma sätt.

% Author's note (verification and reproduction): see the block above the
% hit-list \input at the end of this chapter.  The query table is
% transcribed from the session's query records and the per-provider hit
% counts recorded there.
\begin{table}[p]
  \begin{sidecaption}{Sökfrågor och antal träffar per databas, SRP,
    3 september 2026, verktyget \texttt{scholar}.  S = stödfråga,
    M = motfråga.  Databaser: OpenAlex (OA; inte
    \foreignlanguage{english}{open access}), DBLP, IEEE Xplore (IEEE),
    Web of Science (WoS), Scopus.  Högst 40 träffar hämtades per databas
    och fråga, så 40 betyder \enquote{minst 40}.  Frågan visas i den
    booleska form som skickades till OpenAlex och IEEE Xplore; Web of
    Science fick samma sträng som \texttt{TS=(\dots)}, Scopus som
    \texttt{TITLE-ABS-KEY(\dots)} och DBLP en kort ordlista med samma
    begrepp.  Rader märkta \textsuperscript{b} är samma begrepp ställda
    en gång till i titelformad gestalt, eftersom DBLP bara indexerar
    titlar.  Semantic Scholar avvisade sin nyckel och saknas som kolumn:
    databasen var nere för hela sessionen.
    \textsuperscript{\dag}~Frågan gjordes om (S5r, S7r); den gamla formen
    var för bred och gav mest kemi och medicin.}[tab:sok-queries-srp]
    \centering\footnotesize
    \begin{tabular}{@{}lp{0.44\linewidth}rrrrr@{}}
      \toprule
      & Nyckelord & OA & DBLP & IEEE & WoS & Scopus \\
      \midrule
      S1 & \enquote{single responsibility principle}
         & 40 & 1 & 9 & 14 & 30 \\
      S2 & \enquote{SOLID principles} \textsc{and} \enquote{object-oriented}
         & 38 & 0 & 13 & 12 & 26 \\
      S2\textsuperscript{b} & SOLID principles
         & --- & 28 & --- & --- & --- \\
      S3 & (\enquote{class cohesion} \textsc{or} \enquote{lack of cohesion in methods} \textsc{or} \enquote{cohesion metrics}) \textsc{and} (maintainability \textsc{or} fault \textsc{or} defect)
         & 40 & 0 & 29 & 40 & 40 \\
      S3\textsuperscript{b} & class cohesion metric
         & --- & 22 & --- & --- & --- \\
      S4 & \enquote{responsibility-driven design} \textsc{and} (object-oriented \textsc{or} software)
         & 40 & 3 & 6 & 8 & 21 \\
      S5r\textsuperscript{\dag} & (\enquote{SOLID principles} \textsc{or} \enquote{single responsibility principle} \textsc{or} \enquote{object-oriented design principles}) \textsc{and} (\enquote{software engineering education} \textsc{or} \enquote{computing education} \textsc{or} curriculum \textsc{or} \enquote{teaching programming} \textsc{or} students)
         & 40 & 0 & 9 & 15 & 32 \\
      S6 & (\enquote{information hiding} \textsc{or} \enquote{module decomposition} \textsc{or} \enquote{modular decomposition}) \textsc{and} (criteria \textsc{or} modularity \textsc{or} \enquote{software design})
         & 40 & 0 & 40 & 40 & 40 \\
      S7r\textsuperscript{\dag} & (\enquote{systematic literature review} \textsc{or} \enquote{systematic mapping study} \textsc{or} \enquote{tertiary study}) \textsc{and} (\enquote{SOLID principles} \textsc{or} \enquote{object-oriented design principles} \textsc{or} \enquote{software design principles} \textsc{or} \enquote{design smells} \textsc{or} \enquote{source code metrics})
         & 40 & 1 & 2 & 12 & 15 \\
      S8 & \enquote{single responsibility principle} \textsc{and} (refactoring \textsc{or} \enquote{extract method} \textsc{or} decomposition \textsc{or} modularity \textsc{or} cohesion)
         & 38 & 1 & 4 & 7 & 12 \\
      \midrule
      M1 & (\enquote{single responsibility principle} \textsc{or} \enquote{SOLID principles}) \textsc{and} (criticism \textsc{or} critique \textsc{or} vague \textsc{or} ambiguous \textsc{or} \enquote{ill-defined} \textsc{or} subjective \textsc{or} \enquote{not well defined})
         & 40 & 0 & 0 & 5 & 5 \\
      M2 & (\enquote{SOLID principles} \textsc{or} \enquote{single responsibility principle}) \textsc{and} (empirical \textsc{or} experiment \textsc{or} \enquote{controlled experiment} \textsc{or} evaluation \textsc{or} evidence \textsc{or} \enquote{no effect} \textsc{or} \enquote{negative results} \textsc{or} violation)
         & 40 & 0 & 9 & 18 & 40 \\
      M3 & (\enquote{class cohesion} \textsc{or} \enquote{cohesion metrics} \textsc{or} \enquote{lack of cohesion in methods}) \textsc{and} (criticism \textsc{or} critique \textsc{or} \enquote{no correlation} \textsc{or} \enquote{not correlated} \textsc{or} limitations \textsc{or} \enquote{construct validity} \textsc{or} \enquote{theoretical validation})
         & 40 & 0 & 12 & 36 & 40 \\
      M3\textsuperscript{b} & cohesion metrics evaluation
         & --- & 5 & --- & --- & --- \\
      M4 & (software \textsc{and} (\enquote{design principles} \textsc{or} \enquote{best practices} \textsc{or} \enquote{design heuristics})) \textsc{and} (folklore \textsc{or} \enquote{conventional wisdom} \textsc{or} \enquote{lack of evidence} \textsc{or} unsubstantiated \textsc{or} \enquote{evidence-based software engineering} \textsc{or} anecdotal)
         & 40 & 0 & 26 & 10 & 33 \\
      M5 & (\enquote{code smell} \textsc{or} \enquote{god class} \textsc{or} \enquote{design smell} \textsc{or} \enquote{blob}) \textsc{and} (subjective \textsc{or} \enquote{developer perception} \textsc{or} disagreement \textsc{or} \enquote{inter-rater} \textsc{or} agreement \textsc{or} \enquote{human evaluation})
         & 40 & 0 & 40 & 40 & 40 \\
      M6 & (\enquote{god class} \textsc{or} \enquote{large class} \textsc{or} \enquote{blob} \textsc{or} \enquote{brain class}) \textsc{and} (harmful \textsc{or} \enquote{fault-prone} \textsc{or} \enquote{change-prone} \textsc{or} defects \textsc{or} maintenance)
         & 40 & 0 & 40 & 40 & 40 \\
      \bottomrule
    \end{tabular}
  \end{sidecaption}
\end{table}

Sammanlagt gav sökningarna drygt 1\,400 unika poster.  Alla klassades:
först med en språkmodell mot en beskrivning av frågan och fyra kategorier
--- stöder påståendet, kvalificerar eller motsäger det, angränsande
delämne, annat ämne --- och sedan för hand, där varje post som modellen
satt i den andra kategorin och varje post med låg konfidens lästes.  De
källor som citeras nedan lästes i fulltext, med tre undantag som anges i
\cref{sec:srp-diskussion}.

\section{Resultat}
\label{sec:srp-resultat}

\paragraph{Ursprung.}  Principen härrör från Martin, som formulerar den
ordagrant i kapitel~8 av \emph{Agile Software Development, Principles,
Patterns, and Practices}: \textcquote[ch.~8]{Martin2003Agile}{There should
never be more than one reason for a class to change}.  I samma kapitel
definieras det begrepp påståendet vilar på:
\textcquote[ch.~8]{Martin2003Agile}{In the context of the Single
Responsibility Principle (SRP) we define a responsibility to be
\mkbibquote{a reason for change.}}  Tillskrivningen bekräftas oberoende i
en granskad tidskrift: \textcite{Ardalani2024Supporting} inleder med att
\enquote{according to the single responsibility principle, a module should
have a single responsibility to avoid having more than one reason to
change (Martin, 2003)}.

\paragraph{En rättelse.}  Den artikel som brukar anges som SOLID:s
ursprung, Martins \emph{Design Principles and Design Patterns} från 2000,
innehåller \emph{inte} SRP\autocite{Martin2000Design}.  Dess avsnitt om
klassdesign ger fyra principer --- OCP, LSP, DIP och ISP --- och uttrycket
\enquote{single responsibility} förekommer inte i texten; samma kontroll
av Martins kolumn \emph{Granularity} ger samma resultat.  Den tidigaste
spårbara publiceringen är alltså bokkapitlet, och det är värt att
notera vad det innebär: SRP:s primärkälla är en fackbok, inte en
granskad publikation.

\paragraph{Ursprunget före ursprunget.}  Martin gör inte anspråk på idén
som ny.  I samma kapitel skriver han att principen beskrevs i Tom DeMarcos
och Meilir Page-Jones arbete, där den kallades
kohesion\autocite[ch.~8]{Martin2003Agile}, med fotnoter till deras böcker
\autocite{DeMarco1979Structured,PageJones1988Practical}.  De två böckerna
står här enbart för att dokumentera den tillskrivning Martin själv gör;
de är inte lästa.  Den föregångare som \emph{går} att verifiera är Parnas,
som skriver att \textcquote[s.~1056]{Parnas1972Criteria}{in this context
\mkbibquote{module} is considered to be a responsibility assignment rather
than a subprogram}, och att varje modul i hans föredragna uppdelning
kännetecknas av att den döljer ett designbeslut för de
övriga\autocite[s.~1056]{Parnas1972Criteria}.  Parnas ger alltså modulen
som ansvarstilldelning och uppdelning efter vad som ändras --- men inte
formuleringen \enquote{ett, och endast ett, skäl att ändras}.

\paragraph{Hur etablerad?}  Fyra oberoende slags belägg pekar åt samma
håll.  I forskningen är SRP ett mål som verktyg byggs för:
\textcite{Ardalani2024Supporting} föreslår i \emph{Empirical Software
Engineering} en metod för att dela upp långa metoder i metoder med ett
enda ansvar.  I industrin rapporterar
\textcite{Ampatzoglou2019Applying} en fallstudie där SRP-omstruktureringar
tillämpas på två system med över 1\,500 klasser och kallar SRP
\enquote{a fundamental software engineering principle}, och
\textcite{Anquetil2019Decomposing} beskriver en fastlagd process för att
dela upp gudsklasser hos Siemens, vars planeringssteg ordagrant går ut på
att identifiera de olika ansvar som ska brytas ut.  I undervisningen
rapporterar \textcite{Cabezas2020Using} att SRP lärs ut vid namn i
grundkurser i programmering vid ett ackrediterat program, och
\textcite{Gopal2024Teaching} bekräftar det oberoende från ett annat
lärosäte, med en pilotstudie av \emph{hur} SOLID bäst undervisas.  Och
det underliggande måttet har stöd: \textcite{Iftikhar2024Tertiary}
sammanfattar i en översikt av femton översikter att kodrader, koppling,
komplexitet och kohesionsmåtten ur Chidamber och Kemerers uppsättning är
goda indikatorer på underhållbarhet.

\paragraph{Kritik.}  Motsökningen bär kapitlet, och den gav fyra granskade
invändningar.  Den första är begreppslig och kommer från industristudien
själv: \textcite{Anquetil2019Decomposing} skriver att ansvar i en
gudsklass sällan är klart avgränsade, och att
\enquote{separate responsibilities} i sig är ett subjektivt begrepp.  Den
andra mäter det: \textcite{Alkharabsheh2021Analysing} undersökte hur väl
bedömare och verktyg är överens om att en klass är en gudsklass --- det
kanoniska SRP-brottet --- och fann att överensstämmelsen i allmänhet
saknas, och där den finns är svag.  \Textcite{Hall2014Some} rapporterar
oberoende samma sak för verktyg, för människor och för verktyg mot
människor.  Den tredje invändningen träffar själva förutsägelsen:
\textcite{Olbrich2010Areall} fann att gudsklasser, när effekten
normaliserades för storlek, ändrades mindre och hade färre defekter än
andra klasser --- alltså att de inte nödvändigtvis är skadliga.  Den
fjärde gäller storleksordningen: \textcite{Hall2014Some} rapporterar i
\emph{TOSEM} att alla effekter av de designlukter de undersökt ligger
under tio procent och att andra faktorer betyder mer för fel än
lukterna gör.  Här måste man vara noggrann med räckvidden: de fem lukter
de studerar innefattar inte gudsklassen, så resultatet kvalificerar det
allmänna påståendet att strukturella designregler minskar fel --- det är
inte ett direkt test av SRP.

Till detta kommer två saker som inte är kritik utifrån.  Den ena är att
\textcite{Iftikhar2024Tertiary} skär åt båda hållen: samma översikt som
stöder kohesionsmått noterar att vissa kohesionsmått saknar konsekvent
stöd, och att \textsc{lcom} har pekats ut som teoretiskt ogiltigt som
kohesionsmått.  Den andra, och den tyngsta, är Martins egen: 2014
konstaterar han att formuleringen väcker frågan vad som egentligen
\emph{definierar} ett skäl att ändras, och flyttar kriteriet från koden
till människorna --- det är människor som begär
ändringar\autocite{Martin2014SRP}.  Den omarbetningen bar han vidare in i
\emph{Clean Architecture}\autocite{Martin2017Clean}.

\section{Diskussion och begränsningar}
\label{sec:srp-diskussion}

Ursprunget är entydigt; etableringen är bred men av ett annat slag än
påståendet först antyder.  SRP är etablerad som \emph{praktik} --- den
tillämpas, den undervisas, det byggs verktyg för den --- men sökningen
hittade inget som visar att efterlevnad av SRP i sig ger mätbart bättre
kod.  Varje kvantitativt resultat ovan gäller ombud för SRP: gudsklasser,
kodlukter, kohesionsmått.  Påståendet ska därför formuleras som
\enquote{etablerad praktik}, inte som \enquote{belagd nytta}.

Sex saker begränsar svaret.  Semantic Scholar var nere hela sessionen, och
det är den databas som täcker konferens- och gråzonslitteratur bäst ---
den lucka som betyder mest för en princip vars primärkälla är en fackbok.
Tre källor är lästa bara i sammanfattning: industristudien om
SRP\autocite{Ampatzoglou2019Applying}, studien av gudsklassers
skadlighet\autocite{Olbrich2010Areall} och rapporten om SRP i
grundutbildningen\autocite{Cabezas2020Using}.  Bokens sidnumrering är
inte kontrollerad: kapitlet lästes i en arkiverad utkastversion vars
numrering skiljer sig från den tryckta boken, varför bilagan hänvisar
till kapitel och aldrig till sida.  \emph{Clean Architecture} finns bara
som innehållsförteckning här, så omarbetningen citeras från Martins egen
artikel och inte från boken.  Om SWEBOK eller ACM:s och IEEE:s
kursplaner nämner SOLID vid namn gick inte att fastställa --- det vore
det starkaste enskilda etableringsbelägget, och det är den viktigaste
öppna punkten.  Och klassningen är modellstödd, vilket gör antalet
stödjande poster till en övre gräns: en artikel som \emph{använder} SRP
räknas generöst som stöd för att principen är etablerad.

Ett fynd förtjänar att nämnas och inte citeras.  Motsökningen hittade en
text som direkt bestrider formuleringens riktighet, men den är ett
självdeponerat förtryck utan granskning, utan angiven institution och utan
citeringar.  Att den finns hör till en ärlig redovisning; att låta den stå
som etablerad kritik vore fel.

\section{Slutsats}
\label{sec:srp-slutsats}

Ja på båda frågorna, med förbehåll av två olika slag.  SRP härrör från
\textcite{Martin2003Agile} --- inte, som ofta antas, från hans artikel
från 2000, som saknar principen --- och Martin tillskriver själv
substansen äldre arbete om kohesion och knyter den senare till Parnas.
Principen är etablerad som praktik: den tillämpas i industrin, den är mål
för verktygsbyggande i ledande empiriska tidskrifter, och den undervisas
vid namn i grundkurser.  Men vad som räknas som \emph{ett} ansvar är en
bedömningsfråga, och det är inte en invändning som materialet får tiga om:
begreppet kallas subjektivt i industristudien, oberoende bedömare är
mätbart oense om det kanoniska brottet mot principen, och principens
upphovsman har själv skrivit om formuleringen därför att den inte säger
\emph{vems} skäl att ändra som räknas.  Modulen använder därför SRP som
en fråga att ställa om en funktion --- \enquote{vad skulle få den här att
behöva ändras?} --- och inte som ett kriterium som avgör saken.  Lägg
märke till att förbehållet kom ur motsökningen, och att det bara kom ur
den när frågan ställdes med fältets egna ord för brottet mot principen.
